% Section 10.9, from lectures/L10/appendix/intro_framing_exercises.md and
% lectures/L10/appendix/c_exercises.md.

\section{Exercises}
\label{c10:sec:exercises}

\begin{aside}
The exercises are laid out like this. Exercise~\ref{c10:ex:byhand} and \ref{c10:ex:cpp} belong to
\secref{c10:sec:framing}: the first builds frames by hand, on paper, and the second turns the same
frame into a working C++ implementation. Exercise~\ref{c10:ex:ascan} bridges that frame format and
CAN. Exercise~\ref{c10:ex:stuffing} to \ref{c10:ex:frame} consolidate \secref{c10:sec:can} to
\ref{c10:sec:bittiming}; they are done on paper and build up to assembling a whole frame by hand.
Exercise~\ref{c10:ex:package} writes the \code{can_def} package and then works against it at the
keyboard.
\end{aside}

% ------------------------------------------------------------------------------------------
\exerciseset{Frames by hand and in C++}

\exercise{Constructing two frames by hand}{Reasoning}
\label{c10:ex:byhand}
Start from the frame structure from \secref{c10:sec:framing}:

\begin{textdrawing}
Field  Size      Description
------ --------- -------------------------------------------
SOF    2 bytes   0xA5F7
LEN    1 byte    Payload length
TYPE   1 byte    Message type
DST    1 byte    Destination address
SRC    1 byte    Source address
SEQ    2 bytes   Sequence number
DATA   N bytes   Payload
CHK    2 bytes   Checksum (sum of the preceding bytes)
\end{textdrawing}

The checksum is the sum of every byte from SOF through the last payload byte, as a
\code{std::uint16_t} (so it wraps on overflow). Fields wider than one byte (SOF, SEQ, CHK) are
big-endian. The frame types are \code{Ping = 0x00}, \code{Pong = 0x01}, \code{StatusReq = 0x02} and
\code{StatusResp = 0x03}.

A device at address \code{0x17} sends a status request to a sensor at address \code{0x25}, with
sequence number \code{0x7F05}. That gives \code{TYPE = 0x02}, \code{DST = 0x25}, \code{SRC = 0x17},
\code{SEQ = 0x7F05}, and an empty \code{DATA}.

\xpart{a} Draw the corresponding \code{StatusReq} frame on paper, and fill in every field: SOF, LEN,
TYPE, DST, SRC, SEQ and CHK. Compute the checksum by hand.

\xpart{b} The sensor replies with a 16-bit sensor value, \code{0x3201}, carried in the payload. Draw
the corresponding \code{StatusResp} frame on paper. The rules: TYPE changes to \code{StatusResp}, DST
and SRC swap places, SEQ is unchanged, and DATA is \code{0x3201} (2 bytes). Compute this frame's
checksum by hand as well.

\exercise{The frame in C++}{C++}
\label{c10:ex:cpp}
Implement the frame in C++ as a struct \code{comm::Frame}. Your implementation is validated by an
existing set of unit tests under \file{lectures/L10/appendix/code/}.

\task{1. Examine the file structure}
\begin{itemize}
  \item \file{test/frame_test.cpp} contains the unit tests validating your frame implementation. You
    do not write them; you make them pass.
  \item \file{include/comm/def.hpp} contains the frame fields' offsets (index positions), the field
    sizes and more. The enumeration class \code{comm::FrameType} is to be implemented here.
  \item \file{include/comm/frame.hpp} contains the declaration of the struct \code{comm::Frame}.
  \item \file{source/comm/frame.cpp} is to contain the implementation details of the struct
    \code{comm::Frame}, that is, its method definitions.
\end{itemize}

\task{2. Create the frame types}
Implement the enumeration class \code{comm::FrameType} in \file{comm/def.hpp}, giving each of the
four types the id listed in Exercise~\ref{c10:ex:byhand}. Add a final enumerator \code{Unknown}, one
step beyond the last real type, so that an out-of-range type can be rejected during deserialisation
with a single comparison.

\task{3. Create the frame struct}
Add the member variables of the struct \code{comm::Frame} in \file{comm/frame.hpp}, where
\code{@todo} marks the place: \code{std::uint8_t data[MaxDataLen]}, \code{std::uint16_t seq},
\code{std::uint8_t dst}, \code{std::uint8_t src}, \code{std::uint8_t len} and \code{FrameType type}.

Value-initialise each of them with braces: the numeric members with \code{\{\}}, and \code{type} with
\code{\{FrameType::Unknown\}}, so that a default-constructed frame never holds a type that happens to
be valid.

The widths are not a matter of taste. The unit tests compare a member against a literal of the
field's own width through a template deriving \textbf{one} type from both arguments, so a \code{len}
declared as \code{std::size_t} or a \code{seq} declared as \code{std::uint32_t} does not merely waste
space; the tests stop compiling.

\code{serialize()} and \code{deserialize()} are \textbf{already declared} in that file, documented
and correctly qualified (\code{const noexcept} and \code{noexcept} respectively). Leave those
declarations alone; you write their definitions in steps 4 and 5.

\task{4. Implement the serialisation method}
\Secref{c10:sec:framing}'s \code{checksum16()}, \code{writeToBuf()} and \code{readFromBuf()} are
\textbf{not} part of the code handed out, and both methods below need all three. Copy them into
\file{comm/frame.cpp}, in an anonymous namespace above the method definitions, so that they stay
private to that file:

\begin{cppcode}
namespace
{
// checksum16(), writeToBuf() and readFromBuf() from the framing introduction go here.
} // namespace
\end{cppcode}

Two things about that copying, both of which the compiler only tells you about after the fact:
\begin{itemize}
  \item \code{writeToBuf()} and \code{readFromBuf()} are constrained with \code{std::is_unsigned},
    and nothing in \file{frame.cpp}'s include chain pulls in \code{<type_traits>}. Add that include
    at the top of the file.
  \item Keep \secref{c10:sec:framing}'s warning in mind when you call them: \code{writeToBuf()}
    writes \code{sizeof(T)} bytes derived from the \emph{argument}, not from the field, so cast to
    the field's own width at every call site.
\end{itemize}

Then implement the method \code{comm::Frame::serialize()} so that it writes the following into the
given buffer: \code{SOF = 0xA5F7} (big-endian), LEN, TYPE, DST, SRC, SEQ (big-endian), DATA in the
same order as it sits in the frame's payload buffer, and CHK, computed over all the preceding fields,
byte by byte.

Return the frame's total length if serialisation succeeded, and \code{0U} on failure, which means a
\code{buf} that is null or a \code{bufLen} smaller than the frame to be written. Both have a test of
their own; neither is optional.

\task{5. Implement the deserialisation method}
In the same file, implement the method \code{comm::Frame::deserialize()} so that it validates the
given data, \textbf{in this order}, since every check makes the next one safe to perform:
\begin{itemize}
  \item \code{buf} is not null, and \code{bufLen} is at least \code{MinFrameLen}. Nothing has been
    read yet at this point, so this is what makes reading the header permissible at all. Skip it and
    the reads of SOF and LEN below run off the end of a too-short buffer before any other check gets
    a chance.
  \item SOF is correct.
  \item LEN is within bounds, that is, at most \code{MaxDataLen}.
  \item The given buffer holds the \textbf{whole} frame the header claims. This is a second, stricter
    length check, and the first does not replace it: a buffer truncated after the payload still
    leaves \emph{something} where CHK should be, and reading that as the checksum makes a truncated
    frame look complete.
  \item The frame type is valid, that is, the received byte is below \code{Unknown}.
  \item The checksum matches.
\end{itemize}

Store the given data in the frame, but only once all of it is valid.

Two of those points are checks the tests punish hardest. A missing null check does not make a test
fail, it \textbf{crashes the whole run}, so you see a segfault and no test name. A buffer length
check that stops at the header fails \code{DeserializeRejectsShortBuffer} and
\code{DeserializeRejectsTruncatedPayload} and nothing else, which is a considerably friendlier way of
being told.

\task{6. Validate your implementation}
The unit tests are behind a macro, so nothing runs until you turn it on. Add \code{-DL01} to
\code{CXX_FLAGS} in the makefile, where a \code{Todo} comment marks the place:

\begin{makecode}
CXX_FLAGS := -Wall -Werror -std=c++17 -Iinclude -DL01
\end{makecode}

Then build and run them from a Linux terminal:

\begin{shell}
cd lectures/L10/appendix/code
make
\end{shell}

Until you do, the program reports that the tests are disabled and exits successfully, whether or not
your implementation works.

\begin{aside}
\textbf{If the build stops immediately with this, it is not your code:}

\begin{console}
make[1]: *** No rule to make target 'lib'.  Stop.
make: *** [Makefile:26: ../../../../libs/test/libqacademy_test.a] Error 2
\end{console}

The test framework is a submodule in \file{libs/test}, and a clone without
\code{--recurse-submodules} leaves the directory empty. That it nonetheless \emph{exists} is why the
error is about a missing make target and not about a missing directory: it is the framework's
makefile that has not been fetched yet. Fetch it once, from the repository root, with
\code{git submodule update --init}. The exercise also needs a C++ compiler.
\end{aside}

The tests check both methods against the frames you built by hand in Exercise~\ref{c10:ex:byhand}:
the PING bytes from \secref{c10:sec:framing} and both frames from that exercise, byte by byte; a
round trip through serialisation and deserialisation that has to preserve every field; and the
rejection cases (wrong SOF, wrong checksum, invalid type, truncated buffer, null pointer).

Every passing test is printed in green; a failing one is printed in red and names the expression,
both values and the line, so that you can compare directly against your own paper calculation. Expect
to see all of them fail at first, and go green a group at a time as you work through steps 2--5.

Two things worth knowing in advance:
\begin{itemize}
  \item The tests read \code{frame.type}, \code{frame.dst}, \code{frame.src}, \code{frame.seq},
    \code{frame.len} and \code{frame.data} directly, so those member names have to match. If you have
    named yours something else you are free to change the tests to match.
  \item \code{DeserializeLeavesFrameUntouchedOnFailure} gives \code{deserialize()} a corrupt buffer
    and then checks that the frame's fields are \emph{unchanged}. A \code{deserialize()} copying
    fields into the frame as it parses them passes every other test and fails this one.
\end{itemize}

% ------------------------------------------------------------------------------------------
\exerciseset{From the frame format to CAN}

\exercise{The same frame, as CAN}{Reasoning}
\label{c10:ex:ascan}
Take the \code{StatusResp} frame you built by hand in Exercise~\ref{c10:ex:byhand}\,b), and send it
as a CAN frame instead. Its fields, as a reminder: \code{TYPE = StatusResp}, \code{DST = 0x17},
\code{SRC = 0x25}, \code{SEQ = 0x7F05}, payload 2 bytes with the value \code{0x3201}.

Start from the convention in \secref{c10:sec:broadcast}: the identifier combines a per-type base with
the id of the sensor node the message concerns, that is, the node being asked in a request and the
node replying in a response (node 3 is asked on \code{0x303} and replies on \code{0x403}):

\[ \mathit{ID} = \mathit{TYPE\_BASE} + \mathit{nodeId} \]

The type bases are \code{StatusReq = 0x300}, \code{StatusResp = 0x400}, \code{Ping = 0x500} and
\code{Pong = 0x600}.

\xpart{a} Write out the corresponding CAN frame in the form \code{ID | DLC | DATA}.

\xpart{b} State which fields of \secref{c10:sec:framing}'s frame have no counterpart in the CAN
frame. For every missing field, say what CAN does instead, or explain why CAN needs no equivalent
field.

\xpart{c} Suppose two nodes start transmitting on the CAN bus at exactly the same moment. Explain
briefly how the bus decides which node carries on transmitting, and how the losing node detects that
it has lost. This is the arbitration mechanism from \secref{c10:sec:arbitration}; restate it in your
own words.

\xpart{d} Explain why there is no separate checksum field to fill in when constructing a CAN frame,
unlike \secref{c10:sec:framing}'s frame format.

% ------------------------------------------------------------------------------------------
\exerciseset{Bit stuffing, arbitration and timing}

\exercise{Bit stuffing by hand}{Reasoning}
\label{c10:ex:stuffing}
\xpart{a} Apply the bit stuffing rule to the following 10-bit sequence:

\begin{textdrawing}
0011111111
\end{textdrawing}

Write the complete transmitted bit stream, with all the original bits and every inserted stuff bit
clearly marked.

\xpart{b} A receiver observes the following 9-bit transmitted stream:

\begin{textdrawing}
101111101
\end{textdrawing}

The receiver knows that bit stuffing is in use. Destuff the stream: identify any stuff bits, remove
them, and recover the original sequence of real data bits.

\xpart{c} A receiver observes six dominant (\code{0}) bits in a row at a place where a stuff bit
should have broken a run of five. Explain what the receiver concludes, and what it should do next.

\xpart{d} \Secref{c10:sec:stuffing} points out that the run counter is not cleared at field
boundaries. Suppose a field ends with three identical bits and the next begins with two more of the
same value. Apply the rule and say where the stuff bit lands, counted from the field boundary, and
explain why neither field on its own would have caused it.

\exercise{Decoding a real bit stream}{Reasoning}
\label{c10:ex:decode}
A receiver observes the following 21-bit stream:

\begin{textdrawing}
0 00000010101 0 00 0010 11
\end{textdrawing}

The stream is already destuffed, so that this exercise keeps the focus on the fields' division rather
than on the stuffing. The gaps are only there for readability; the real CAN bus has no gaps between
the fields.

\xpart{a} Split the stream into its named fields, with the widths from \secref{c10:sec:frameformat}:
SOF, ID, RTR, IDE, r0, DLC. Then determine the identifier in hexadecimal, the DLC value, and which
field has begun in the two bits left over once all the named fields are accounted for.

\xpart{b} From the DLC in \textbf{a)}, determine how many data bits the receiver expects in total,
and how many further data bits remain before the CRC field begins.

\xpart{c} Suppose the DLC had instead been \code{0000}. Which field would then follow immediately
after the control field?

\exercise{Arbitration between three nodes}{Reasoning}
\label{c10:ex:arb3}
Three nodes try to transmit at the same SOF: node R with ID \code{0x155}, node S with ID
\code{0x154}, and node T with ID \code{0x1D5}.

\xpart{a} Work through the arbitration bit by bit, in the same style as
\secref{c10:sec:arbitration}'s worked example. Determine which node wins, and at which identifier bit
each losing node drops out.

\xpart{b} The two losing nodes drop out at different bit positions. Explain why by stating where each
losing node's identifier first differs from the winner's identifier, and why that first differing bit
in particular decides when the node loses arbitration.

\xpart{c} A losing node has to stop driving the bus, and the interesting question is \emph{when}.
(\Chapref{c11:ch} introduces the port pair \code{tx_bus}/\code{bus_en} that the book's controller
stops driving with, and \chapref{c16:ch}--\ref{c17:ch} build the logic; here you reason at the
protocol level.) Explain, in one or two sentences, what each losing node's \code{can_controller} has
to do on the very first clock cycle after it detects the discrepancy, and why that timing matters.

\note[Clue]Think through what ``the moment'' means in \secref{c10:sec:arbitration}'s worked
arbitration example.

\exercise{Reasoning about the sample point}{Reasoning}
\label{c10:ex:sample}
This book samples every bit 70~\% into the bit period, not at 0~\% (the very start of the period) or
100~\% (the very end).

\xpart{a} Explain what could go wrong if a receiver sampled at 0~\%. Start from signal propagation,
why a bit period needs a minimum duration, and whether the transmitted value has had time to settle.

\xpart{b} Explain what could go wrong if a receiver sampled at 100~\%, on the last possible clock
pulse. Make the explanation specific to the fact that in the book's design every node has its own
local bit timer, that the timer is resynchronised only at SOF, and that it is never resynchronised
mid-frame (real CAN resynchronises on edges mid-frame as well; the book's controller does not).

\xpart{c} Arbitration requires every competing node to read the bus during the \emph{same} bit period
as it drives it. Explain why a sample point too near the start of the period is risky, why a sample
point too near the end of the period is risky, and why the danger grows the further into the frame
you get.

\exercise{Constructing a CAN frame by hand}{Reasoning}
\label{c10:ex:frame}
Exercise~\ref{c10:ex:stuffing} to \ref{c10:ex:sample} worked with individual rules. This one
assembles a complete frame from the frame diagram in \secref{c10:sec:frameformat}.

\begin{aside}
\textbf{A remark about the CRC.} You are given the CRC-15 value for every frame below. Computing a
15-bit CRC over some 30-odd bits by hand is mechanical rather than instructive; in \chapref{c13:ch}
you work through the recursion properly, over a 4-bit sequence, when the engine itself is the
subject.
\end{aside}

\xpart{a} Take the CAN frame you produced in Exercise~\ref{c10:ex:ascan}: \code{ID = 0x425},
\code{DLC = 2}, data \code{32 01}. Write out every field as bits, in transmission order, with the
widths from \secref{c10:sec:frameformat}: SOF; ID (11 bits, MSB first); RTR; IDE and r0; DLC; data
(MSB first, byte 0 first); CRC (given: \code{010101100100011}); CRC delimiter, ACK slot, ACK delimiter
and EOF.

\xpart{b} How many bits is that in total, before any stuffing?

\xpart{c} Now apply the bit stuffing rule from Exercise~\ref{c10:ex:stuffing} over SOF through the
CRC field, remembering that the run counter is not cleared at field boundaries. How many stuff bits
are inserted, and how many bits does the frame occupy on the wire?

\xpart{d} At the book's bit rate of 1~Mbit/s, how long does it take to send that frame?

\xpart{e} Repeat \textbf{a)} to \textbf{d)} for a frame \textbf{with no payload}: \code{ID = 0x100},
\code{DLC = 0}, CRC \code{011100000001010}. Note which parts of the frame are structurally unchanged
despite the missing data field, and where the CRC field now begins.

\xpart{f} Compare your two answers to \textbf{c)}. The second frame carries no data at all, and yet
its stuffed region needs \emph{more} stuff bits than the first. Explain why, in terms of what the bit
stuffing rule actually reacts to.

% ------------------------------------------------------------------------------------------
\exerciseset{Working with the package}

\exercise{\code{can_def} and its derivations}{VHDL}
\label{c10:ex:package}
Every constant in \code{can_def} can be traced back to something this chapter defined. These
questions write the package, check that the derivations hold when the inputs change, and extend it
with values \secref{c10:sec:frameformat} named but the package does not yet hold.

\xpart{a} Write \file{controller/can_def.vhd}, with the constants and subtypes
\secref{c10:sec:candef} sets out. If you were at the session you already have it; if you were not,
this is where it comes from, and every chapter from \chapref{c11:ch} onwards assumes it exists.

This is the only module in the book handed out in full rather than as a specification, and for a
reason worth mentioning: a package contains no behaviour to design, so there is nothing between the
values and the code for you to work out. What is new here is the VHDL itself, since a package
declaration is a construct the book has not used before. Write it out by hand rather than pasting it
in. The names are the next eight chapters' vocabulary, and writing them once is cheaper than looking
them up eight times.

\xpart{b} Suppose the DE0-CV's clock were 40~MHz instead of 50~MHz, with the bit rate still
1~Mbit/s. Recompute \code{TICKS_PER_BIT} and \code{SAMPLE_TICK} (still 70~\%).

\xpart{c} Suppose instead that the clock stays at 50~MHz but the bus runs at 125~kbit/s, a real CAN
``low speed'' rate. Recompute both constants.

\xpart{d} \code{TICKS_PER_BIT} is \code{CLOCK_FREQ_HZ / BIT_RATE_HZ}, with VHDL's integer division
for \code{natural}, which truncates. Give a pair of clock and bit rate where the division does
\emph{not} come out even, and explain in one sentence why that would be a problem for this design,
even though the VHDL would still compile and elaborate without complaint.

\xpart{e} Add a constant \code{BIT_PERIOD_NS}, the length of a CAN bit period in nanoseconds, derived
from \code{BIT_RATE_HZ} rather than hardcoded. Then state the relationship between it,
\code{TICKS_PER_BIT} and the 50 MHz system clock's period of 20~ns, and confirm that the three agree.

\xpart{f} \Secref{c10:sec:frameformat} states the CRC delimiter, the ACK field and EOF as 1, 2 and 7
bits respectively, but the package holds none of them: \code{can_controller} will use those numbers
as bare literals. Add \code{CRC_DELIM_BITS}, \code{ACK_BITS} and \code{EOF_BITS}, and say which of the
two forms a later reader is better served by.

\xpart{g} Check that the package actually compiles. Install GHDL if you have not already
(appendix~\ref{app:sim}), and then run \code{make build-project} from the root of the group's
repository. You should see:

\begin{console}
--> controller/can_def.vhd analyzes cleanly
\end{console}

followed by the two \file{bridge/} files handed out, and after that every testbench reported as
skipped, since the modules they check do not exist yet:

\begin{console}
0 testbench(es) run, 8 skipped.
\end{console}

Now break something deliberately, a missing semicolon or a misspelled type, rerun, and read what GHDL
says about it. Put it back.

That is the whole verification story for a package: a package has no behaviour to simulate, so ``it
analyses'' is all there is to check. \Chapref{c11:ch} introduces the tool properly, and from there
every module gets a testbench too.
